\PartQuestion{Which biological mechanisms motivate which computational abstractions, and where does each mapping stop?}
Part III moves from biology to typed software. It begins with real regulatory dynamics, where history determines abundance and some parameters cannot be identified from protein data alone. It then builds generated structure, evolutionary search, typed genomes, operational semantics, resource ledgers and intervention-based evidence. Each chapter ends with a contract that later chapters can test. The recurring lesson is that a biological name supplies a question, not an equivalence: a software gate is not a promoter, and a credit score is not a mechanism.
